% Esquema de modelagem completo: constelação, cinco estrelas e matriz de barramento.
% Compilar com: latexmk -xelatex modelagem.tex
\documentclass[11pt]{article}
\usepackage{fontspec}
\usepackage[brazil]{babel}
\usepackage[a4paper,landscape,margin=1.6cm]{geometry}
\usepackage{amsmath}
\usepackage{tikz}
\usetikzlibrary{shapes.multipart,positioning,calc}
\usepackage{booktabs}
\usepackage{enumitem}
\usepackage{underscore}
\pagestyle{empty}
\setlist{itemsep=1pt,topsep=3pt}
\newcommand{\pagina}[2]{\noindent{\Large\bfseries #1}\par\vspace{2pt}\noindent{\small #2}\par\vspace{8pt}}
% Notação dos esquemas (a mesma do artigo da disciplina): cada tabela é uma caixa com
% o nome no topo e os atributos embaixo; (PK) marca a chave primária artificial e (FK)
% as chaves estrangeiras da tabela de fatos. A tabela de fatos tem borda mais grossa.
\tikzset{
  dim/.style={rectangle split, rectangle split parts=2, draw, font=\scriptsize,
              text width=3.95cm, inner sep=2.4pt, rectangle split part align={center,left},
              rectangle split part fill={black!7,white}},
  fato/.style={dim, line width=1.1pt, text width=4.35cm, rectangle split part fill={black!17,white}},
  nome/.style={draw, font=\scriptsize\bfseries, minimum width=2.55cm, minimum height=0.66cm,
               align=center, fill=black!7, inner sep=2pt},
  nomefato/.style={nome, line width=1.1pt, minimum width=2.75cm, fill=black!17},
  liga/.style={line width=0.55pt},
  rollup/.style={dashed, black!55, line width=0.6pt},
  papel/.style={font=\tiny\itshape, fill=white, inner sep=1pt, text=black},
}
% ponto k de n ao longo do lado esquerdo / direito da tabela de fatos
\newcommand{\ladoesq}[3]{($(#1.north west)!{#2/(#3+1)}!(#1.south west)$)}
\newcommand{\ladodir}[3]{($(#1.north east)!{#2/(#3+1)}!(#1.south east)$)}

\begin{document}

\pagina{Constelação de fatos}{Rede Municipal Paulistana: oferta escolar e demanda não atendida · SCC0245 · Grupo 7. Cinco tabelas de fatos e catorze dimensões. Linhas contínuas: fato e dimensão; as duas linhas entre Matrícula e Data são os dois papéis (início e fim). Tracejado: rollup encolhido conformado.}
\begin{center}\resizebox{\linewidth}{!}{\begin{tikzpicture}
% tabelas de fatos (linha central)
\node[nomefato] (FM) at (-7.5,0) {Tabela de Fatos\\Matrícula};
\node[nomefato] (FD) at (-2.5,0) {Tabela de Fatos\\Desfecho};
\node[nomefato] (FT) at (1.5,0) {Tabela de Fatos\\Ocupação de Turma};
\node[nomefato] (FF) at (6.0,0) {Tabela de Fatos\\Fila};
\node[nomefato] (FP) at (10.5,0) {Tabela de Fatos\\População};
% dimensões de cima
\node[nome] (DA) at (-9.6,3.2) {Dimensão Data\\(dois papéis)};
\node[nome] (AL) at (-5.6,3.2) {Dimensão\\Ano Letivo};
\node[nome] (ES) at (-2.5,3.2) {Dimensão\\Escola};
\node[nome] (SE) at (0.6,3.2) {Dimensão\\Série};
\node[nome] (ET) at (4.3,3.2) {Dimensão\\Etapa};
\node[nome] (TE) at (8.25,3.2) {Dimensão\\Tempo};
% dimensões de baixo
\node[nome] (ID) at (-7.5,-3.4) {Dimensão\\Idade};
\node[nome] (TU) at (-2.6,-3.4) {Dimensão\\Turno};
\node[nome] (TR) at (2.0,-3.4) {Dimensão\\Território};
\node[nome] (SC) at (6.3,-3.4) {Dimensão Setor\\Censitário};
\node[nome] (FE) at (10.5,-3.4) {Dimensão\\Faixa Etária};
% bugigangas do fato de matrícula
\node[nome] (PA) at (-11.2,1.4) {Dimensão\\Perfil Aluno};
\node[nome] (DE) at (-11.2,0) {Dimensão\\Deficiência};
\node[nome] (SM) at (-11.2,-1.4) {Dimensão Situação\\Matrícula};

% fato -> dimensão
\draw[liga] ([xshift=-5pt]FM.north) -- ([xshift=8pt]DA.south);
\draw[liga] ([xshift=1pt]FM.north) -- ([xshift=16pt]DA.south);
\foreach \d in {AL,ES,SE,TU,TR,ID,PA,DE,SM} \draw[liga] (FM) -- (\d);
\foreach \d in {AL,ES,SE,TR} \draw[liga] (FD) -- (\d);
\foreach \d in {AL,ES,SE,TU,TR} \draw[liga] (FT) -- (\d);
\foreach \d in {ET,TE,TR} \draw[liga] (FF) -- (\d);
\foreach \d in {TE,SC,FE,TR} \draw[liga] (FP) -- (\d);

% rollups encolhidos conformados (tracejado)
\draw[rollup] (DA) -- node[papel, above=1pt]{rollup} (AL);
\draw[rollup] (DA.north) |- node[papel, pos=0.75]{rollup: dia $\leq$ mês} (8.25,4.55) -- (TE.north);
\draw[rollup] (SE) -- node[papel, above=1pt]{rollup} (ET);
\draw[rollup] (SC) -- node[papel, above=1pt]{rollup} (TR);
\draw[rollup] (ID.south) |- node[papel, pos=0.75]{rollup: idade $\leq$ faixa IBGE} (10.5,-5.0) -- (FE.south);
\end{tikzpicture}
}\end{center}

\clearpage
\pagina{Esquema estrela Matrícula}{Grão: um vínculo aluno-turma em um ano letivo. Fatos: quantidade de matrículas (campo artificial = 1, SOMA) e dias de matrícula (SOMA).}
\begin{center}\resizebox{!}{0.78\textheight}{\begin{tikzpicture}[node distance=0.32cm]
\node[fato] (F) at (0,0) {\textbf{Tabela de Fatos Matrícula}
\nodepart{two}(FK) chave ano letivo\\(FK) chave data início\\(FK) chave data fim\\(FK) chave escola\\(FK) chave território\\(FK) chave série\\(FK) chave turno\\(FK) chave idade\\(FK) chave perfil aluno\\(FK) chave deficiência\\(FK) chave situação\\id turma \hfill\textit{(degenerada)}\\código matrícula \hfill\textit{(degenerada)}\\quantidade matrículas \hfill\textit{(artificial)}\\dias de matrícula};

\node[dim, anchor=north east] (A) at (-3.35,6.15) {\textbf{Dimensão Ano Letivo}\nodepart{two}(PK) chave ano letivo\\ano\\qualidade código escola\\observação};
\node[dim, below=of A] (D) {\textbf{Dimensão Data}\nodepart{two}(PK) chave data\\data, dia, dia da semana, nome do dia\\mês, nome mês\\trimestre, semestre, ano\\fim de semana};
\node[dim, below=of D] (E) {\textbf{Dimensão Escola}\nodepart{two}(PK) chave escola\\código escola, código sintético\\versão, ano início, ano fim, versão atual\\nome escola, tipo escola, rede\\distrito, DRE, setor SME, situação};
\node[dim, below=of E] (T) {\textbf{Dimensão Território}\nodepart{two}(PK) chave território\\distrito\\DRE, participação DRE\\município, UF};
\node[dim, below=of T] (S) {\textbf{Dimensão Série}\nodepart{two}(PK) chave série\\modalidade, segmento, série\\ciclo, ordem, educação especial\\código etapa, etapa, nível};

\node[dim, anchor=north west] (TU) at (3.35,6.15) {\textbf{Dimensão Turno}\nodepart{two}(PK) chave turno\\turno\\jornada\\origem turno};
\node[dim, below=of TU] (I) {\textbf{Dimensão Idade}\nodepart{two}(PK) chave idade\\idade\\faixa escolar, grupo etário\\faixa IBGE};
\node[dim, below=of I] (P) {\textbf{Dimensão Perfil Aluno}\nodepart{two}(PK) chave perfil aluno\\código combinação\\sexo, raça e cor\\altas habilidades, idade estimada};
\node[dim, below=of P] (DF) {\textbf{Dimensão Deficiência}\nodepart{two}(PK) chave deficiência\\código combinação\\possui deficiência, quantidade\\17 indicadores (autismo, ...)};
\node[dim, below=of DF] (SI) {\textbf{Dimensão Situação Matrícula}\nodepart{two}(PK) chave situação\\código combinação\\situação matrícula, desfecho\\12 indicadores (aprovado, ...)};

\draw[liga] (A.east) -- \ladoesq{F}{1}{7};
\draw[liga] ([yshift=4pt]D.east) -- node[papel, pos=0.35]{início} \ladoesq{F}{2}{7};
\draw[liga] ([yshift=-4pt]D.east) -- node[papel, pos=0.35]{fim} \ladoesq{F}{3}{7};
\draw[liga] (E.east) -- \ladoesq{F}{4}{7};
\draw[liga] (T.east) -- \ladoesq{F}{5}{7};
\draw[liga] (S.east) -- \ladoesq{F}{6}{7};
\foreach \n/\k in {TU/1,I/2,P/3,DF/4,SI/5} \draw[liga] (\n.west) -- \ladodir{F}{\k}{5};
\end{tikzpicture}
}\end{center}
\noindent\small Hierarquias: dia $\leq$ mês $\leq$ trimestre $\leq$ semestre $\leq$ ano; distrito $\leq$ DRE $\leq$ município $\leq$ UF; escola $\leq$ tipo de unidade $\leq$ rede; série $\leq$ etapa $\leq$ nível e série $\leq$ ciclo; turno $\leq$ jornada; idade $\leq$ faixa escolar, idade $\leq$ grupo etário e idade $\leq$ faixa IBGE.

\clearpage
\pagina{Esquema estrela Fila por vaga}{Grão: distrito, mês de referência e etapa. Todos os fatos são estoques (semiaditivos): SOMA em Território e Etapa, valor de um único mês no Tempo.}
\begin{center}\resizebox{0.9\linewidth}{!}{\begin{tikzpicture}[node distance=0.45cm]
\node[fato] (F) at (0,0) {\textbf{Tabela de Fatos Fila}
\nodepart{two}(FK) chave tempo\\(FK) chave território\\(FK) chave etapa\\formato planilha \hfill\textit{(degenerada)}\\matrículas\\matrículas em processo\\demanda total\\demanda sem opção\\demanda sem opção até 30 dias\\demanda sem opção mais de 30 dias\\demanda com opção por unidade};
\node[dim, anchor=north east] (TE) at (-3.35,2.55) {\textbf{Dimensão Tempo}\nodepart{two}(PK) chave tempo\\data referência, ano mês\\mês, nome mês\\bimestre, trimestre, semestre\\ano\\fila publicada, observação};
\node[dim, below=of TE] (TR) {\textbf{Dimensão Território}\nodepart{two}(PK) chave território\\distrito\\DRE, participação DRE\\município, UF};
\node[dim, anchor=west] (ET) at (3.35,0) {\textbf{Dimensão Etapa}\nodepart{two}(PK) chave etapa\\código\\etapa\\nível\\ordem};
\draw[liga] (TE.east) -- \ladoesq{F}{1}{3};
\draw[liga] (TR.east) -- \ladoesq{F}{2}{3};
\draw[liga] (ET.west) -- (F.east);
\end{tikzpicture}
}\end{center}
\noindent\small Hierarquias: mês $\leq$ trimestre $\leq$ semestre $\leq$ ano e mês $\leq$ bimestre $\leq$ ano; distrito $\leq$ DRE $\leq$ município $\leq$ UF; etapa $\leq$ nível.

\clearpage
\pagina{Esquema estrela Desfecho da matrícula}{Grão: escola, ano letivo e série (fato agregado da Matrícula). Contagens: SOMA. Taxas: não aditivas, recalculadas pelas contagens em cada nível.}
\begin{center}\resizebox{0.9\linewidth}{!}{\begin{tikzpicture}[node distance=0.45cm]
\node[fato] (F) at (0,0) {\textbf{Tabela de Fatos Desfecho}
\nodepart{two}(FK) chave ano letivo\\(FK) chave escola\\(FK) chave território\\(FK) chave série\\matrículas\\aprovados, reprovados, abandonos\\transferidos, concluintes, falecidos\\outras saídas, em curso, sem desfecho\\\textit{taxa aprovação (derivada)}\\\textit{taxa reprovação (derivada)}\\\textit{taxa abandono (derivada)}};
\node[dim, anchor=north east] (A) at (-3.35,2.35) {\textbf{Dimensão Ano Letivo}\nodepart{two}(PK) chave ano letivo\\ano\\qualidade código escola\\observação};
\node[dim, below=of A] (E) {\textbf{Dimensão Escola}\nodepart{two}(PK) chave escola\\código escola, versão\\nome escola, tipo escola, rede\\distrito, DRE, situação\\...};
\node[dim, anchor=north west] (T) at (3.35,2.35) {\textbf{Dimensão Território}\nodepart{two}(PK) chave território\\distrito\\DRE, participação DRE\\município, UF};
\node[dim, below=of T] (S) {\textbf{Dimensão Série}\nodepart{two}(PK) chave série\\modalidade, segmento, série\\ciclo, ordem, educação especial\\código etapa, etapa, nível};
\draw[liga] (A.east) -- \ladoesq{F}{1}{2};
\draw[liga] (E.east) -- \ladoesq{F}{2}{2};
\draw[liga] (T.west) -- \ladodir{F}{1}{2};
\draw[liga] (S.west) -- \ladodir{F}{2}{2};
\end{tikzpicture}
}\end{center}

\clearpage
\pagina{Esquema estrela Ocupação de Turma}{Grão: uma turma em um ano letivo (fato agregado da Matrícula). Contagens: SOMA. Alunos por turma: não aditivo, soma dos alunos ativos dividida pela soma das turmas ativas.}
\begin{center}\resizebox{0.9\linewidth}{!}{\begin{tikzpicture}[node distance=0.45cm]
\node[fato] (F) at (0,0) {\textbf{Tabela de Fatos Ocupação de Turma}
\nodepart{two}(FK) chave ano letivo\\(FK) chave escola\\(FK) chave território\\(FK) chave série\\(FK) chave turno\\id turma \hfill\textit{(degenerada)}\\quantidade turmas \hfill\textit{(artificial)}\\matrículas no ano\\matrículas ativas em junho\\turmas ativas em junho\\\textit{alunos por turma (derivado)}};
\node[dim, anchor=north east] (A) at (-3.35,3.25) {\textbf{Dimensão Ano Letivo}\nodepart{two}(PK) chave ano letivo\\ano\\qualidade código escola\\observação};
\node[dim, below=of A] (E) {\textbf{Dimensão Escola}\nodepart{two}(PK) chave escola\\código escola, versão\\nome escola, tipo escola, rede\\distrito, DRE, situação\\...};
\node[dim, below=of E] (T) {\textbf{Dimensão Território}\nodepart{two}(PK) chave território\\distrito\\DRE, participação DRE\\município, UF};
\node[dim, anchor=north west] (S) at (3.35,2.35) {\textbf{Dimensão Série}\nodepart{two}(PK) chave série\\modalidade, segmento, série\\ciclo, educação especial\\código etapa, etapa, nível};
\node[dim, below=of S] (TU) {\textbf{Dimensão Turno}\nodepart{two}(PK) chave turno\\turno\\jornada\\origem turno};
\draw[liga] (A.east) -- \ladoesq{F}{1}{3};
\draw[liga] (E.east) -- \ladoesq{F}{2}{3};
\draw[liga] (T.east) -- \ladoesq{F}{3}{3};
\draw[liga] (S.west) -- \ladodir{F}{1}{2};
\draw[liga] (TU.west) -- \ladodir{F}{2}{2};
\end{tikzpicture}
}\end{center}

\clearpage
\pagina{Esquema estrela População residente}{Grão: setor censitário, faixa etária e data de referência do Censo 2022. Estoque (semiaditivo): SOMA em setor, distrito e faixa; um único Censo no tempo.}
\begin{center}\resizebox{0.9\linewidth}{!}{\begin{tikzpicture}[node distance=0.45cm]
\node[fato] (F) at (0,0) {\textbf{Tabela de Fatos População}
\nodepart{two}(FK) chave tempo\\(FK) chave setor censitário\\(FK) chave território\\(FK) chave faixa etária\\quantidade de pessoas};
\node[dim, anchor=north east] (TE) at (-3.35,2.75) {\textbf{Dimensão Tempo}\nodepart{two}(PK) chave tempo\\data referência, ano mês\\mês, nome mês\\bimestre, trimestre, semestre\\ano\\fila publicada, observação};
\node[dim, below=of TE] (SC) {\textbf{Dimensão Setor Censitário}\nodepart{two}(PK) chave setor censitário\\código setor\\código distrito IBGE\\nome distrito IBGE\\distrito\\município, UF};
\node[dim, anchor=north west] (TR) at (3.35,2.75) {\textbf{Dimensão Território}\nodepart{two}(PK) chave território\\distrito\\DRE, participação DRE\\município, UF};
\node[dim, below=of TR] (FE) {\textbf{Dimensão Faixa Etária}\nodepart{two}(PK) chave faixa etária\\faixa IBGE\\idade inicial, idade final\\variável IBGE\\ordem};
\draw[liga] (TE.east) -- \ladoesq{F}{1}{2};
\draw[liga] (SC.east) -- \ladoesq{F}{2}{2};
\draw[liga] (TR.west) -- \ladodir{F}{1}{2};
\draw[liga] (FE.west) -- \ladodir{F}{2}{2};
\end{tikzpicture}
}\end{center}
\noindent\small Hierarquias: setor censitário $\leq$ distrito $\leq$ DRE $\leq$ município $\leq$ UF; a faixa etária é o topo de idade $\leq$ faixa IBGE.

\clearpage
\pagina{Matriz de barramento}{Dimensões conformadas usadas por cada tabela de fatos.}
\begin{center}\large
\begin{tabular}{lccccc}
\toprule
\textbf{Dimensão} & \textbf{Matrícula} & \textbf{Fila} & \textbf{Desfecho} & \textbf{Turma} & \textbf{População}\\\midrule
Território & $\times$ & $\times$ & $\times$ & $\times$ & $\times$ \\
Ano Letivo (ano) & $\times$ & & $\times$ & $\times$ & \\
Tempo (mês) & & $\times$ & & & $\times$ \\
Data (dia, dois papéis) & $\times$ & & & & \\
Escola & $\times$ & & $\times$ & $\times$ & \\
Série & $\times$ & & $\times$ & $\times$ & \\
Etapa (rollup de Série) & & $\times$ & & & \\
Turno & $\times$ & & & $\times$ & \\
Idade & $\times$ & & & & \\
Faixa Etária (rollup de Idade) & & & & & $\times$ \\
Setor Censitário & & & & & $\times$ \\
Perfil do Aluno, Deficiência, Situação & $\times$ & & & & \\
\bottomrule
\end{tabular}
\end{center}
\end{document}
